\section{Wafer-Scale Phylogeny Trials}


To conduct these simulations, we used an extension of the island-model genetic algorithm framework presented in \cite{moreno2024trackable}.
This framework organizes simulated organisms into independent subpopulations hosted on each PE.
Synchronous tournament selection of size 2 is applied within each PE.
Between generations, agents are migrated between PEs asynchronously.

%The purpose of this simulation was to generate large-scale data using a simple model that can be explicitly manipulated to explore evolutionary regimes differing substantially in the depth of shared ancestry between organisms (known as phylogenetic richness).

A simple genome model was used, comprised of a single floating point value representing fitness.
Reproduction occurred asexually, with this value subject to a normally-distributed mutation with 30\% probability.
(In the case of the purifying-regime treatment, the negative absolute value of a sampled mutation delta was applied.)
In tournament competitions to select parents for the subsequent generation, the parent was selected according to the higher fitness value with 5\% probability and was otherwise selected randomly.
Genomes were augmented with hstrat annotations comprising 64 single-bit markers, managed at runtime using a \texttt{tilted} curation policy, which favors dense retention of more recent marker data \cite{moreno2024algorithms}.

To generate data representative of differing evolutionary conditions, two treatments were applied: an \textit{adaptive regime} and a \textit{purifying regime}.
Under the adaptive regime, mutations increasing fitness were allowed, introducing the possibility of selective sweeps.
By contrast, the purifying regime treatment allowed only deleterious mutations.
Previous work with this system demonstrated the purifying regime to exhibit substantially greater phylogenetic richness, in line with expectations for selective sweeps to diminish population-level phylogenetic diversity \cite{moreno2024trackable}.

Reported neutral conditions correspond to purifying regime.

Materials are available via the Open Science Framework at\url{https://osf.io/yn3wx} \cite{foster2017open}.
% To enable the simulation scale necessary for this work, we employed Graphics Processing Unit (GPU) and Wafer-Scale Engine (WSE) hardware accelerators, which contain large collections of worker processors specialized for numerical computation.
WSE hardware was Cerebras CS-3 in Cerebras Cloud with Cerebras SDK v1.4.0 \cite{buitrago2021neocortex}.

Further methodological details on WSE phylogeny reconstructions are available as \cite{singhvi2025scalable}.
